% ===================================================================
% Section 2. Setting
% Sources: research/theory/T1 §1 (and §2.3 scope remark, §5 side conditions),
%          T2 §1 (and §2.1-2.2, §4.1, §7), T7 §1, T3 §1 (contexts preview only).
% Proofs of the lemmas: sections/app-setting.tex.
% ===================================================================
\providecommand{\ACCEPT}{\textsf{ACCEPT}}
\providecommand{\REJECT}{\textsf{REJECT}}
\providecommand{\ESCALATE}{\textsf{ESCALATE}}
\providecommand{\Steps}{\mathrm{Steps}}
\providecommand{\vars}{\mathrm{vars}}
\providecommand{\Ttrust}{T_{\mathrm{trust}}}
\providecommand{\Fres}[1]{F^{(#1)}_{\mathrm{res}}}

\section{Setting: judgments, steps, rules, provers, and three channels}\label{sec:setting}

This section fixes the vocabulary of the paper. Nothing in it is deep; the point
is a set-up in which the later theorems are short. Three questions are kept apart
from the start: \emph{what} is learned (a step set, a consequence relation, or a
set of valuations; \Cref{sec:setting:hyp}), \emph{from what} (three channels of
different logical form; \Cref{sec:setting:channels}), and \emph{against whom} (an
adaptive prover that searches for wrongly accepted steps;
\Cref{sec:setting:protocol}). Two running examples, propositional natural
deduction and school algebra (\Cref{sec:setting:practice}), recur throughout.

% -------------------------------------------------------------------
\subsection{Judgments, steps and closure}\label{sec:setting:steps}

$\Fm$ is the set of formulas over atoms $p,q,r,p_0,p_1,\dots$ with
$\neg,\wedge,\vee,\to,\bot,\top$; first-order languages are used for arithmetic
and physics. A \emph{context} $\Gamma$ is a finite set of formulas.

\begin{definition}[Judgments, positions]\src{T1 \S1.1; T2 \S1}\label{def:setting:judgment}
A \emph{judgment} is $\Gamma\der\varphi$ (single conclusion) or $\Gamma\der\Delta$
(multiple conclusion, $\Delta$ finite); $J$ denotes a countable set of judgments.
Contexts are part of judgments, so hypothetical reasoning, and ``assume the air
pressure is $0$'', lives inside the step language. A \emph{position}
$\pos{\Gamma}{\Delta}$ asserts all of $\Gamma$ and denies all of $\Delta$.
Following \citet{restall2005multiple}, $\Gamma\der\Delta$ says that \emph{the
position $\pos{\Gamma}{\Delta}$ is out of bounds}; a position is \emph{coherent}
if it is in bounds. When $\bot$ is present it counts as always denied, so
$\pos{\Gamma}{\emptyset}$ is coherent iff $\Gamma\nder\bot$. Without $\bot$,
``coherent'' means \emph{non-trivial} ($\Gamma$ does not entail every formula;
for a structural relation, equivalently $\Gamma\nder q$ for an atom $q$ not in
$\Gamma$).
\end{definition}

Restall's reading gives inference data and coherence data opposite polarity: an
inference says a position is \emph{out of} bounds, a coherence datum certifies
that one is \emph{in} bounds. This is the bilateral picture of
\citet{smiley1996rejection} and \citet{rumfitt2000yes} and the multiple-conclusion
logic of \citet{shoesmith1978multiple}, anticipated in Carnap's ``full
formalization'' \citep{carnap1943formalization}. \Cref{sec:coherence} shows that
this polarity decides Carnap's categoricity problem for a learner.

\begin{definition}[Steps, rule sets, closure]\src{T1 \S1.1}\label{def:setting:step}
A \emph{step} $P\step j$ has a finite set $P\subseteq J$ of premises and a
conclusion $j\in J$; $S$ is the set of steps, and a \emph{rule set} (step set) is
any $R\subseteq S$. An \emph{$R$-derivation} $\pi$ of $j$ from $B\subseteq J$ is a
finite tree with root $j$ whose nodes are leaves in $B$ or conclusions of steps
$P\step j'\in R$ whose premises $P$ label their children ($P=\emptyset$ allowed);
$\Steps(\pi)$ is the set of steps used. The \emph{closure} $\Cl_R(B)$ is the least
$X\supseteq B$ such that $P\step j\in R$ and $P\subseteq X$ imply $j\in X$; the
\emph{derived steps} are $\Sound(R)=\{P\step j:\ j\in\Cl_R(P)\}$. The
\emph{target} is an unknown rule set $\Rstar$: its members are the \emph{valid}
steps, and the members of $\Sound(\Rstar)$ the \emph{derivable} ones.
\end{definition}

\begin{lemma}[Closure facts; trivial]\src{T1 \S1.1}\label{lem:setting:closure}
\textup{(a)} $\Cl_R(B)$ is the set of judgments with an $R$-derivation from $B$;
$\Cl_R$ is extensive, monotone, idempotent and finitary, and monotone in $R$.
\textup{(b)} $R\subseteq\Sound(R)$ and $\Cl_{\Sound(R)}=\Cl_R$; hence $\Sound(R)$
is the largest step set $A$ with $\Cl_A=\Cl_R$.
\end{lemma}
Proof in Appendix~\ref{app:setting}.

Two formats recur. In the \emph{formula format} (Hilbert style) $J=\Fm$ and a step
is $\varphi_1,\dots,\varphi_k\step\psi$. In the \emph{sequent format} (natural
deduction in sequent style; \citealt{gentzen1935untersuchungen,prawitz1965natural})
judgments are sequents and discharge is an ordinary step, e.g.\ $\to$I is
$\{\Gamma,\varphi\der\psi\}\step\Gamma\der\varphi\to\psi$. A rule set induces a
relation on sets of formulas: $A\der_R\varphi$ iff $\varphi\in\Cl_R(A)$ (formula
format), or iff $(A_0\der\varphi)\in\Cl_R(\emptyset)$ for a finite
$A_0\subseteq A$ (sequent format). In the formula format $\der_R$ is a finitary
consequence relation by \Cref{lem:setting:closure}(a): arguments chain, so cut is
built in; for the natural-deduction calculi below, weakening and cut are
admissible and the same holds. $R$ is \emph{$A$-coherent} if $A\nder_R\bot$.

\paragraph{Reading.} \Cref{sec:search} shows that a verifier is safe under
derivations of every length iff everything it accepts lies in $\Sound(\Rstar)$ (T1
Lemma~1.1). So ``learning which inferences are valid'' has a sharp target: a set
between $\Rstar$ and $\Sound(\Rstar)$.

% -------------------------------------------------------------------
\subsection{Schemas and anti-unification}\label{sec:setting:schemas}

Steps are encoded as ground terms over a first-order signature $\Omega$, e.g.\
$\mathsf{s}(\mathsf{and}(p_0,p_1),p_0)$ for an $\wedge$E step; object-language
atoms are \emph{constants} of $\Omega$.

\begin{definition}[Schemas]\src{T1 \S\S1.3,\,5; T7 \S1.1}\label{def:setting:schema}
A \emph{schema} $\sigma$ is a term with metavariables $x$, $y$, $A$, $B$, \dots, and
$\inst(\sigma)$ is its set of ground instances. $\sigma\succeq\tau$ ($\sigma$ is
at least as general) iff $\tau=\sigma\theta$ for a substitution $\theta$;
$\sigma\equiv\tau$ iff both (equality up to renaming). $\sigma$ is \emph{pure} if
no atom constant occurs in it; its \emph{purification} $\sigma^\circ$ replaces
distinct atoms by distinct fresh metavariables. A \emph{guarded schema}
$(\sigma,\Psi)$ carries guards $\Psi\subseteq\Phi$, a fixed finite family of
decidable predicates on ground steps (freshness; ``the facts entail $a\neq0$''),
and $\inst(\sigma,\Psi)$ is the set of instances of $\sigma$ satisfying every
guard in $\Psi$. A
\emph{calculus} is a finite set $\Sigma$ of (guarded) schemas, with
$R_\Sigma=\bigcup_{\sigma\in\Sigma}\inst(\sigma)$; a \emph{tagged} step $(i,s)$
records the schema $i$ that it cites.
\end{definition}

Anti-unification is the learning algorithm of this paper. For a tuple
$\bar u=(u_1,\dots,u_n)$ of ground terms, $\mathcal A(\bar u)=f(\mathcal A(\bar
u^1),\dots,\mathcal A(\bar u^r))$ if all $u_j$ have root $f$ of arity $r$
($\bar u^k$ is the tuple of $k$-th arguments); otherwise $\mathcal A(\bar
u)=z_{\bar u}$, a metavariable indexed by the tuple, so equal tuples get equal
metavariables.

\begin{theorem}[Least general generalization; \citealt{plotkin1970note,reynolds1970transformational}]\src{T1 \S1.3}\ \status{known}\label{thm:setting:lgg}
Every nonempty finite set $P=\{t_1,\dots,t_n\}$ of ground terms has a least
general generalization $\lgg(P)$, a schema with $P\subseteq\inst(\lgg(P))$ and
$\lgg(P)\preceq\tau$ whenever $P\subseteq\inst(\tau)$. It is unique up to
renaming and computed by $\mathcal A(t_1,\dots,t_n)$. Hence
$\inst(\lgg(P))=\bigcap\{\inst(\tau):P\subseteq\inst(\tau)\}$: anti-unification is
the \emph{closure algorithm} for single-schema instance sets.
\end{theorem}

For a schema $t$ let $\mu(t):=|t|-|\vars(t)|$, where $|t|$ counts symbol
occurrences (metavariables included) and $\vars(t)$ is the set of distinct
metavariables; so $\mu\ge0$, $\mu(x)=0$, and $\mu(t)=|t|$ for ground $t$.

\begin{lemma}[Rank]\src{T1 Lemma 1.2}\label{lem:setting:rank}
If $\tau=\sigma\theta$ then $\mu(\tau)\ge\mu(\sigma)$, with equality iff $\theta$
restricted to $\vars(\sigma)$ is a renaming. Hence every strictly increasing
generalization chain $\tau=\sigma_0\prec\sigma_1\prec\dots\prec\sigma_m$ has
$m\le\mu(\tau)-\mu(\sigma_m)\le\mu(\tau)$.
\end{lemma}

\begin{lemma}[When the lgg recovers the schema]\src{T1 Lemma 1.3}\label{lem:setting:recover}
Let $t_j=\sigma\theta_j$ $(1\le j\le n)$ be ground instances of $\sigma$. Then
$\lgg(t_1,\dots,t_n)\preceq\sigma$, with $\equiv$ iff both \emph{witness events}
hold: \textup{(R)} for each metavariable $x$ of $\sigma$, the roots of
$\theta_1x,\dots,\theta_nx$ are not all equal; \textup{(D)} for each $x\neq y$
there is $j$ with $\theta_jx\neq\theta_jy$.
\end{lemma}

\noindent\emph{Proof idea.} Substituting $\theta x$ for the occurrences of $x$
adds more symbols than distinct metavariables unless $\theta x$ is a metavariable,
which gives the rank lemma (folklore in the Plotkin--Reynolds lattice).
Anti-unifying the $t_j$ copies $\sigma$ and puts $\mathcal
A(\theta_1x,\dots,\theta_nx)$ at $x$; this is a metavariable iff (R) holds, and two
of them coincide iff (D) fails. Full proofs in Appendix~\ref{app:setting}. By the
rank lemma, ascending generalization chains are finite, so lggs exist for infinite
$P$ as well. The positive-data theory of \Cref{sec:caution,sec:imitation} runs on
counting witness events.

\begin{lemma}[Guarded closure; trivial]\src{T1 \S5, side conditions}\label{lem:setting:guard}
For nonempty finite $P\subseteq S$ let $\Psi_P=\{\phi\in\Phi:\phi\equiv1\text{ on
}P\}$. Then $\inst(\lgg(P),\Psi_P)$ is the least set of the form
$\inst(\sigma,\Psi)$ containing $P$.
\end{lemma}

So a learned guard is the strongest guard satisfied by \emph{every} datum, and one
datum that omits a side condition deletes it (\Cref{ex:setting:alg}). Proof in
Appendix~\ref{app:setting}.

% -------------------------------------------------------------------
\subsection{Hypothesis classes}\label{sec:setting:hyp}

\begin{definition}[Hypothesis classes]\src{T2 \S1, \S4.1; T1 \S1.2, \S3}\label{def:setting:hyp}
Hypotheses come in three guises.
(i)~\emph{Step sets} $\Hc\subseteq2^S$: e.g.\ single schemas
$H_1=\{\inst(\sigma)\}\cup\{\emptyset\}$, tagged calculi
$H^{\mathrm{tag}}_k=\{\bigcup_{i\le k}\{i\}\times\inst(\sigma_i)\}$, untagged
unions $H_k=\{\bigcup_{l\le k}\inst(\tau_l)\}$, finite unstructured classes. A
learned step scorer (a process reward model) is a step set, and need not be
closed under chaining.
(ii)~\emph{Consequence relations} $\der$ (finitary, reflexive, monotone,
cut-closed), equivalently finitary closure operators $C$; $\der$ is
\emph{structural} if $\Gamma\der\varphi$ implies $s\Gamma\der s\varphi$ for every
substitution $s$ of formulas for atoms \citep{los1958remarks,wojcicki1988theory}.
Each $R$ induces $\der_R$; many $R$ induce the same relation.
(iii)~\emph{Meanings as sets of valuations}: a set $V$ of admissible valuations
$v:\Fm\to\{0,1\}$ (not necessarily compositional), with $\Gamma\models^m_V\Delta$
iff no $v\in V$ has $v[\Gamma]=1$ and $v[\Delta]=0$; dually, $\Val(S')$ is the
set of valuations violating no sequent in $S'$.
The \emph{version space} of data $D$ is $\VS(D)=\{h\in\Hc: h\text{ consistent
with }D\}$ ($=\{h:D\subseteq h\}$ for positive data). $\Hc$ is \emph{realizable}
if it contains the target; $w$ is a prior and $\wstar=w(\Rstar)$.
\end{definition}

\emph{Structurality is uniformity.} If $\sigma$ is pure, $\inst(\sigma)$ is closed
under substitutions $s$ for atoms, since the atoms of $\sigma\theta$ come from
$\theta$: $s(\sigma\theta)=\sigma\theta'$ with $\theta'x=s(\theta x)$. So a learner
that generalizes human steps to pure schemas is structural. A schema that mentions
an atom is not, and can be unsound without trivializing: added to a complete
calculus for $\CPC$, $x\vee p_0\step x$ derives $\neg p_0$ but stays inside
$\Cn(\neg p_0)$ (T1 Prop.~2.3, scope remark). Purification restores the link
(\Cref{sec:search}).

\emph{Inference data are negative data about valuations.} The sequent
$\Gamma\step\Delta$ excludes exactly the valuations realizing
$\pos{\Gamma}{\Delta}$; a coherence certification or a world report is
\emph{positive} data about valuations (T2 \S4.1). Imitation and coherence inform
the learner from opposite sides, and \Cref{sec:coherence} makes this precise as
``Carnap's problem is Gold's problem in the dual space''.

% -------------------------------------------------------------------
\subsection{Practice and the running examples}\label{sec:setting:practice}

\begin{definition}[Practice]\src{T7 Defs 1.1--1.2; T1 \S6}\label{def:setting:practice}
The \emph{target} is a finite calculus $\Sigma^*$ (the \emph{genuine rules}), with
$\Rstar=R_{\Sigma^*}$. The \emph{systematic fallacies} are a finite set $F$ of
schemas, $F\cap\Sigma^*=\emptyset$, each with an instance outside
$\Sound(\Rstar)$. The \emph{practice} is $\Sigma^P=\Sigma^*\sqcup F$, with
$R^P=\Rstar\cup R_F$. \emph{Human data} are i.i.d.\ tagged steps: tag
$i\in\Sigma^P$ has probability $\pi_i$, and given $i$ the step is an instance of
$i$ with probability $1-a_i$ and otherwise an arbitrary step outside $\inst(i)$
tagged $i$ (\emph{sporadic noise}). Alternatively, the data are a \emph{text}, an
enumeration of $R^P$ \citep{gold1967language}.
\end{definition}

Fallacy tags are idealized as \emph{latent but distinct}: a student who affirms
the consequent ``cites'' a rule that is in fact affirming the consequent. If
fallacies are cited under a genuine rule's name, the per-tag class becomes $H_k$
(\Cref{sec:imitation,sec:twotier}).

\begin{example}[Propositional natural deduction]\label{ex:setting:nd}
$\Sigma^*$ is a natural-deduction calculus for $\CPC$ ($\wedge$I/E, $\vee$I/E,
$\to$I, $\to$E $=$ modus ponens (MP), $\neg$I/E, ex falso, double negation). In
the formula format MP is $\varphi,\ \varphi\to\psi\step\psi$. A student sometimes
\emph{affirms the consequent}, $\mathrm{AC}:\ \psi,\ \varphi\to\psi\step\varphi$;
so $F=\{\mathrm{AC}\}$. The data are steps such as $p\wedge q\step p$ tagged
$\wedge$E. The two instances $p\wedge q\step p$ and $(r\to p)\wedge p\step r\to p$
already satisfy (R) and (D), so their lgg is $A\wedge B\step A$
(\Cref{lem:setting:recover}). (In the sequent format contexts are sets, and
anti-unification modulo set structure is finitary rather than unitary; we flag
this where it matters.)
\end{example}

\begin{example}[School algebra]\label{ex:setting:alg}
Judgments are equations $s=t$ between terms built from $+$, $-$, $\cdot$, $/$,
powers, $\sqrt{\ }$, numerals and variables, under finite sets of \emph{facts} such as $x\neq0$. Terms
denote partial functions ($a/0$ is undefined), and $s=t$ is true under the facts
iff $s,t$ are Kleene-equal (both undefined, or both defined and equal) at every
real assignment satisfying them. Steps are instances of rewrite schemas such as
$(a+b)^2\step a^2+2ab+b^2$ and $a/a\step1$ with guard $a\neq0$, plus trusted
equational steps (symmetry, transitivity, congruence) and trusted numeral
evaluation. The student commits the \emph{freshman's dream} $(a+b)^2\step
a^2+b^2$, an invalid schema, and \emph{unguarded cancellation} $a/a\step1$, a
genuine rule used without its guard (the instance $x/x\step1$ fails at $x=0$). If
the student cites ``cancellation'' for both uses, one datum such as
$(y+1)/(y+1)\step1$ under no facts drops the guard from the learned rule
(\Cref{lem:setting:guard}); with distinct latent tags the unguarded rule is learned
as a schema of its own. Either way, positive data learn it as a rule.
\end{example}

These are experiments A and B of \Cref{sec:experiments}.

% -------------------------------------------------------------------
\subsection{Three channels}\label{sec:setting:channels}

\begin{definition}[The channels]\src{T1 \S1.2; T2 \S1, \S2.1--2.2; T7 Defs 1.2--1.4; L3 Lemma A}\label{def:setting:channels}
\leavevmode
\begin{enumerate}[label=(\Alph*),leftmargin=*]
\item[(P)] \emph{Practice/imitation} delivers human steps $X_1,X_2,\dots$,
  possibly tagged with rule names, possibly with sporadic noise
  (\Cref{def:setting:practice}). A noise-free datum lies in $R^P$; no datum says
  that a step is invalid.
\item[(C)] \emph{Coherence} supplies a family $\Ac$ of \emph{designated contexts},
  finite premise sets certified coherent for the target ($A\nder_{\Rstar}\bot$),
  or, bilaterally, \emph{designated positions} $\pos{\Gamma}{\Delta}$ certified in
  bounds ($\bot\notin\Cl_{\Rstar}(\Gamma)$ and
  $\Delta\cap\Cl_{\Rstar}(\Gamma)=\emptyset$). By default $\Ac$ is fixed and
  known, independently of the target. A hypothesis $h$ is \emph{$\Ac$-coherent}
  if $A\nder_h\bot$ for all $A\in\Ac$. A \emph{negative bag} is $\Steps(\pi)$ for
  an argument $\pi$ from some $A\in\Ac$ to $\bot$; it contains an invalid step, or
  chaining would give $A\der_{\Rstar}\bot$.
\item[(W)] The \emph{world} is an evaluable fragment $J_W\subseteq J$ with truth
  function $W:J_W\to\{0,1\}$, plus optional \emph{trusted steps} $\Ttrust$. Its
  kinds: (W1) \emph{truth values on a fragment} (closed propositional formulas,
  $\Delta_0$ sentences, a decision procedure); (W2) \emph{counterexample objects},
  structures in which every line of a given argument can be evaluated; (W3)
  \emph{computation with fresh randomness}, e.g.\ testing a polynomial identity at
  a point drawn after the claim is committed; (W4) \emph{simulation}, a validated
  numerical solution of a less idealized model (\Cref{sec:physics}).
\end{enumerate}
\end{definition}

Coherence is enforced \emph{only} on designated contexts: deriving $\bot$ from an
undesignated supposition is a reductio and is never penalized. A negative bag is a
multiple-instance label \citep{dietterich1997solving}, ``not all of these steps
are valid''; it collapses to an ordinary negative example only for a learner that
can test $A\step\bot$ directly against cut-closed hypotheses, so the bag structure
matters exactly for step verifiers. Designations that are data about the target
carry more information (\Cref{sec:coherence}). Objects (W2) turn bags into
singletons: if a structure makes the leaves of an argument true and its conclusion
false, then starting at the conclusion and repeatedly moving to a false premise
ends at a step with true premises and a false conclusion, which alone is blamed
(the \emph{descent} lemma, \Cref{sec:twotier,sec:informal}). Fresh randomness (W3)
is sound against adaptive provers only under conditions on commitment order,
degree bounds and the field (L3 Lemma~A, \Cref{sec:search}); Schwartz--Zippel
bounds its error \citep{schwartz1980fast,zippel1979probabilistic}.

\begin{assumption}[(WS), sound evaluation]\src{T7 \S1.3}\label{asm:setting:WS}
For each designated $\pos{\Gamma}{\Delta}$, the partial evaluation with $e(j)=1$
if $j\in\Gamma$ or $W(j)=1$, and $e(j)=0$ if $j\in\Delta\cup\{\bot\}$ or $W(j)=0$,
is well defined and extends to a total $V:J\to\{0,1\}$ under which no step of
$\Rstar\cup\Ttrust$ has true premises and a false conclusion; and
$\Ttrust\subseteq\Sound(\Rstar)$.
\end{assumption}

In formal mathematics $V$ is truth in the intended structure $\mathfrak M$: (WS)
says that designated positions are true of $\mathfrak M$, $W$ is truth in
$\mathfrak M$ on its fragment, and target and trusted steps are $\mathfrak
M$-sound.

\begin{lemma}[(WS) and coherence; trivial]\src{T7 \S1.3}\label{lem:setting:WS}
Under (WS) every designated position is in bounds for $\Rstar$. Conversely, if
$J_W=\emptyset$, $\Ttrust\subseteq\Sound(\Rstar)$ and every designated position is
in bounds, then (WS) holds.
\end{lemma}
Proof in Appendix~\ref{app:setting}.

\begin{table}[t]
\centering\small
\begin{tabularx}{\textwidth}{@{}lXXX@{}}
\toprule
 & (P) practice & (C) coherence & (W) world \\
\midrule
datum & ``this step is used'' & an argument from a designated context to $\bot$ & a truth value; an object; a random test; a simulation \\
logical form & positive instance & negative bag (``some step is invalid'') & singleton blame with objects; derivation-level otherwise \\
blind to & systematic errors & errors coherent with the target & what lies outside the evaluable fragment \\
\bottomrule
\end{tabularx}
\caption{The three channels (\Cref{def:setting:channels}).}
\label{tab:setting:channels}
\end{table}

\begin{example}[\Cref{ex:setting:nd} continued]\label{ex:setting:nd2}
(C) Designate $A_1=\{q,\ p\to q,\ p\to\bot\}$: AC gives $p$, MP then gives $\bot$.
The bag $\{q,p\to q\step p;\ p,p\to\bot\step\bot\}$ does not say which step is at
fault, and the rival diagnosis ``MP is the fallacy'' reads $\to$ as $\leftarrow$
(\Cref{sec:twotier}). (W2) The valuation $v(p)=0$, $v(q)=1$ makes $A_1$ true and
the line $p$ false, so descent blames AC alone. (W1) The closed instance
$\top,\ \bot\to\top\step\bot$ of AC has true premises and a false conclusion.
\end{example}

\begin{example}[\Cref{ex:setting:alg} continued]\label{ex:setting:alg2}
(C) Designate the empty fact set, with $\bot$ any equation between distinct
numerals. The freshman's dream gives $(1+1)^2=1^2+1^2$, i.e.\ $4=2$ after trusted
evaluation: a singleton bag. Unguarded cancellation gives $2\cdot(0/0)=2\cdot1=2$
and $2\cdot(0/0)=(2\cdot0)/0=0/0=1$, so $1=2$; this bag also contains the genuine
rules $a\cdot(b/c)\step(a\cdot b)/c$, $a\cdot0\step0$ (for defined $a$) and
$a\cdot1\step a$. (W2) An evaluator that observes ``undefined'' blames
$0/0\step1$ alone. (W3) The freshman's dream is off by $2ab$, so a uniformly random
point of $S^2$, $S\subseteq\mathbb Z$ finite, exposes it with probability at least
$1-2/|S|$; over $\mathbb F_2$, where it is an identity, never: the field matters.
The instance $x/x\step1$ fails only at $x=0$, which random evaluation essentially
never hits; the object $x=0$ does. In experiment A, coherence without (W) removes
value errors like the freshman's dream but almost never learns guards that concern
only definedness (such as $a\neq0$ for $0/a\step0$); what survives is sound for a
total reading of division over $\mathbb C$ with $1/0=0$, a coherent alternative
meaning \status{computed; \Cref{sec:experiments}}.
\end{example}

\paragraph{Reading.} The channels are the user's three proposed signals:
imitation of human inferences, a coherence loss, and occasional feedback from the
world. They are not three noisy versions of one signal. (P) can only confirm, (C)
can only condemn sets, and (W) can condemn single steps when it supplies objects.
Each later section says what one of them fixes and what it provably cannot.

% -------------------------------------------------------------------
\subsection{Provers, verifiers, soundness and cost}\label{sec:setting:protocol}

\begin{definition}[Protocol]\src{T1 \S1.2; T2 \S2.2}\label{def:setting:protocol}
In rounds $t=1,2,\dots$ a \emph{prover} submits $q_t\in S$ and the
\emph{verifier} answers $a_t\in\{\ACCEPT,\REJECT,\ESCALATE\}$; on \ESCALATE{} an
oracle returns $\mathbf 1[q_t\in\Rstar]$ (possibly noisy), which joins the
history. The \emph{acceptance region} $\hat R_t$ is the set of steps accepted
outright at time $t$; a reasoner chains them and derives $\Cl_{\hat R_t}(B)$. The
prover may be randomized, adaptive and computationally unbounded; it knows the
target, the practice, the verifier's code and the history, but not the verifier's
future coins, future human data, or the oracle's future (fresh) noise. It models
any proof search, the learner's own included. It is \emph{honest} if all its
queries are valid.
\end{definition}

\begin{definition}[Soundness]\src{T1 Def 1.2; T7 \S1.5}\label{def:setting:soundness}
Fix a class $\Hc$ of targets and, per target, a set $S^\dagger$ of admissible steps
(default $S^\dagger=\Rstar$). A verifier is \emph{$\delta$-sound} if for every
target and every prover
$\Prob[\exists t:\ a_t=\ACCEPT\wedge q_t\notin S^\dagger]\le\delta$ (a
time-uniform guarantee); \emph{deterministically} (or $0$-)\emph{sound} if
$\delta=0$; and \emph{uniformly $\delta$-sound} if some event $G$ with
$\Prob(G)\ge1-\delta$, depending only on the human data and the verifier's coins,
has no prover ever getting a step outside $S^\dagger$ accepted on $G$.
\emph{Truth-soundness} takes $S^\dagger=\Rstar$ for the genuine $\Sigma^*$,
\emph{practice-soundness} $S^\dagger=R^P$.
\end{definition}

Uniform $\delta$-soundness implies $\delta$-soundness, since every prover's bad
event lies outside $G$. Asking for $S^\dagger=\Rstar$ is stronger than the
condition $\hat R_t\subseteq\Sound(\Rstar)$ that chaining needs; our verifiers
meet it. Soundness is relative to the target calculus, not to truth: improving on
$\Rstar$ is the job of (C) and (W). Since coherence is $\Pi_1$ (T2 Thm~5.3), no
computable learner is both eventually complete and sound at all times relative to
the full-depth residue of fallacies (T7 Thm~5.7, \Cref{sec:twotier}); soundness is
then stated relative to a refutation budget.

\begin{definition}[Depth-relative soundness]\src{T7 Defs 1.5--1.6}\label{def:setting:depth}
Assume (WS). A \emph{$d$-refutation} of a finite schema set $B$ is a derivation of
total size $\le d$, relative to a designated position, whose leaves evaluate to
$1$, whose steps lie in $R_B\cup\Ttrust$, and whose conclusion evaluates to $0$
($\bot$, a denied judgment, or a $W$-false one). $B$ is \emph{$d$-clean} if it has
none. The \emph{depth-$d$ residue} $\Fres{d}=\{\tau\in F:\Sigma^*\cup\{\tau\}\text{
is }d\text{-clean}\}$ decreases in $d$. A verifier is \emph{depth-$d$ sound} if it
is (uniformly) $\delta$-sound for $S^\dagger=\Rstar\cup R_{\Fres{d}}$.
\end{definition}

\begin{definition}[Cost]\src{T1 Def 1.3; T2 \S2.2; T7 \S1.4}\label{def:setting:cost}
(a)~\emph{Escalations}: the number of rounds in which a valid query is not
accepted outright (in expectation, for randomized verifiers). $\Esc(\Hc\mid D)$ is
the infimum over $0$-sound verifiers of the supremum over targets in $\Hc$
containing the data $D$, and over honest query sequences, of this number;
$\Esc(\Hc)=\Esc(\Hc\mid\emptyset)$ is the \emph{escalation dimension}. Invalid
queries are not charged.
(b)~\emph{Detections}: rounds in which an argument $\pi$ from some $A\in\Ac$ to
$\bot$ with $\Steps(\pi)\subseteq\hat R_t$ is exhibited; \emph{incompleteness
errors}: rounds presenting a valid step outside $\hat R_t$.
(c)~\emph{Mind changes}: rounds with $\hat R_{t+1}\neq\hat R_t$; those that remove
an accepted step are \emph{retractions} \citep{kelly1996logic}.
(d)~\emph{Samples} $N$ (human steps) and \emph{calls} to a refutation oracle that
returns a $d$-refutation of a queried set or reports none.
\end{definition}

\begin{example}[One bad step suffices]\label{ex:setting:exploit}
If AC is accepted, a prover derives any $\varphi$ from no premises: $\top$ and
$\varphi\to\top$ are theorems, and AC concludes $\varphi$. If unguarded
cancellation is accepted, it derives $1=2$ (\Cref{ex:setting:alg2}). A verifier
whose errors are rare on the human distribution still accepts such a step every
time a prover asks for it. This is Prior's \emph{tonk}
\citep{prior1960runabout} as an attack on a learned checker; \Cref{sec:search}
turns it into theorems.
\end{example}

\paragraph{Reading.} This is how the set-up takes the user's ``compute is not a
constraint, but it must actually work''. The prover is unbounded and adaptive, so
soundness is a worst-case property: reliable learning \citep{rivest1988learning},
perfect selective classification \citep{elyaniv2010foundations} and the ``accept''
side of KWIK \citep{li2008knows}, applied to steps. Its price is paid in
escalations (human time), samples, detections and oracle calls, counted
separately. Average-case accuracy is not enough (\Cref{sec:search}).

% -------------------------------------------------------------------
\subsection{Contexts: a preview}\label{sec:setting:contexts}

Contexts let one argue inside a supposition or an idealization without being
penalized for contradicting the world. Following \citet{mccarthy1993notes}, write
$\ist(c,\varphi)$ for ``$\varphi$ holds in context $c$''; contexts form a tree with
root $@$, the actual or intended situation (T3 \S1). A \emph{supposition}
$\SUP(A)$ refines its parent by $A$, and discharge is exact:
$\ist(\SUP(A),\psi)$ iff $A\to\psi$ holds in the parent (T3 Prop.~1.5); deriving
$\bot$ there is a reductio of $A$, not an incoherence. An \emph{idealization}
$\DEF$ deforms named parameters of a declared model family toward an ideal value
(``air pressure $\to0$''). It may contradict the root, yet does not explode while
it is well posed: truth in it is truth in every admissible completion (a filter of
models). Coherence is enforced on the root and on contexts anchored to it by
certified exports, never on suppositions or idealizations as such; designating
``idealization plus full background'' as coherent would force every structural
learner to give up a classical rule \emph{globally} (T2 Prop.~7.1), so only
consistent chunks are designated \citep{brown2004chunk}. Contexts, export
certificates and the physics checker are the subject of \Cref{sec:physics}.
